\section{Trace-Driven Virtual PIM Evaluation}
\label{sec:method}

\subsection{Problem Framing}

We study virtual PIM offloading during autoregressive LLM decoding. Let a decode trace be a sequence of region-level records
\[
\mathcal{T} = \{r_t\}_{t=1}^{T},
\]
where each record contains a region identity, measured latency, estimated bytes moved, arithmetic intensity, context length, batch size, and auxiliary metadata derived from the runtime. We do not attempt to model a full physical PIM system. Instead, we use these records to evaluate a virtual decision problem: for each decode region, should execution remain on the GPU or be offloaded to a hypothetical PIM backend under a prior-work-grounded latency model?

\subsection{Trace Collection}

Our framework instruments real decoding runs from open-source runtimes and collects region-level records for attention and MLP-dominated portions of each decode step. Each trace record stores:
\begin{itemize}
\item measured latency,
\item estimated floating-point work and bytes moved,
\item current context length and batch size,
\item runtime-derived KV-related proxies such as $\kv$ bytes and memory pressure.
\end{itemize}
This design gives us a runtime-grounded description of decode behavior without requiring proprietary serving stacks or custom hardware counters.

\subsection{Virtual PIM Backend}

The virtual backend is intentionally transparent. It uses a prior-work-grounded cost model rather than a cycle-accurate microarchitecture. The GPU-side assumptions are anchored to a roofline-style A100 reference, while the PIM side is parameterized by bandwidth, compute throughput, and fixed orchestration overheads. This makes the backend suitable for trace-grounded replay, signal comparison, and sensitivity analysis, but not for direct hardware-speedup claims.

For each trace record $r$, the GPU latency estimate is
\[
L_{\mathrm{gpu}}(r)=
\begin{cases}
\ell_{\mathrm{trace}}(r), & \text{if measured region latency is available}, \\
\max\!\left(\frac{\mathrm{FLOPs}(r)}{\Theta_{\mathrm{gpu}}}, \frac{\mathrm{Bytes}(r)}{B_{\mathrm{gpu}}}\right), & \text{otherwise}.
\end{cases}
\]
The virtual PIM latency uses the trace-grounded GPU latency as an anchor and then rescales its compute- and memory-dominant shares:
\[
L_{\mathrm{pim}}(r)=
\max\!\left(
L_{\mathrm{gpu}}(r)\cdot s_c(r)\cdot \frac{\Theta_{\mathrm{gpu}}}{\Theta_{\mathrm{pim}}^{\mathrm{eff}}(r)},
L_{\mathrm{gpu}}(r)\cdot s_m(r)\cdot \frac{B_{\mathrm{gpu}}}{B_{\mathrm{pim}}^{\mathrm{eff}}(r)}
\right)
\;+\;
\tau_{\mathrm{transfer}}
\;+\;
\tau_{\mathrm{sync}},
\]
where $s_c(r)$ and $s_m(r)$ are the compute- and memory-share terms derived from the GPU roofline decomposition. For attention regions, $B_{\mathrm{pim}}^{\mathrm{eff}}(r)$ receives a context-scaled bandwidth bonus; for non-attention and MLP regions, $\Theta_{\mathrm{pim}}^{\mathrm{eff}}(r)$ is reduced by fixed penalties to reflect their weaker fit to memory-centric execution. The default parameters and provenance split are summarized in \Cref{tab:backend-measured-modeled,tab:backend-params}.

\begin{table}[t]
\centering
\caption{Virtual backend inputs grouped by provenance.}
\label{tab:backend-measured-modeled}
\begin{tabular}{p{0.22\linewidth}p{0.37\linewidth}p{0.27\linewidth}}
\toprule
Category & Fields & Source / Interpretation \\
\midrule
Trace-measured & region label, measured latency, context length, batch size & Directly observed in the collected decode traces \\
Trace-derived & bytes moved, FLOPs, arithmetic intensity, estimated KV bytes, KV bytes/token, memory-pressure proxy & Derived from runtime metadata and trace-side estimators \\
Modeled only & roofline compute/memory shares, effective PIM bandwidth/throughput modifiers, transfer/sync overheads, estimated PIM latency & Virtual replay assumptions used only for what-if comparison \\
\bottomrule
\end{tabular}
\end{table}

\begin{table}[t]
\centering
\caption{Default virtual backend parameters, role, and provenance.}
\label{tab:backend-params}
\begin{tabular}{p{0.28\linewidth}p{0.13\linewidth}p{0.24\linewidth}p{0.24\linewidth}}
\toprule
Parameter & Default & Role & Provenance \\
\midrule
GPU peak BW & 1935 GB/s & Roofline anchor for memory time & Prior-work A100-style roofline anchor \\
GPU peak throughput & 312 TFLOP/s & Roofline anchor for compute time & Prior-work A100-style roofline anchor \\
PIM peak BW & 2400 GB/s & Baseline virtual PIM memory rate & Memory-centric default: modestly above GPU BW to encode bandwidth advantage without assuming equal compute \\
PIM peak throughput & 18 TFLOP/s & Baseline virtual PIM compute rate & Conservative low-throughput default to avoid overstating compute-heavy gains \\
Transfer overhead & 0.020 ms & Per-dispatch orchestration cost & Conservative midpoint varied again in preset sensitivity \\
Sync overhead & 0.010 ms & Per-dispatch synchronization cost & Conservative midpoint varied again in preset sensitivity \\
Attention BW bonus & 0.20 & Long-context attention memory modifier & Qualitative prior reflecting stronger KV pressure at long context \\
FC parallelism penalty & 0.25 & MLP parallelism penalty & Deliberate penalty to avoid overstating non-attention and MLP benefit \\
Non-attention penalty & 0.10 & Throughput penalty outside attention & Deliberate penalty to keep non-attention gains conservative \\
\bottomrule
\end{tabular}
\end{table}


This paper treats the virtual backend as a methodological enabler rather than as the main novelty. All central claims are therefore phrased in terms of regime characterization and relative policy behavior, not exact hardware performance.

\subsection{Policies}

We evaluate the following policy families:
\begin{itemize}
\item \textbf{GPU-only}: never offload.
\item \textbf{Static attention}: offload all attention regions.
\item \textbf{Online predictor}: a lightweight dynamic predictor using runtime features.
\item \textbf{Adaptive feature}: a realistic feature-based policy using context, bytes moved, arithmetic intensity, latency, and region bias.
\item \textbf{KV regime}: a realistic policy that augments runtime features with explicit KV-related signals such as estimated $\kv$ bytes, $\kv$ bytes per token, and memory pressure proxy.
\item \textbf{Oracle}: an upper bound that chooses the lower-cost device for each record under the virtual backend.
\end{itemize}

The key comparison in this paper is not between a sophisticated learned scheduler and a naive baseline, but between different kinds of realistic decision signals. We use this comparison to test whether static attention rules are sufficient, whether context length alone is already a strong coarse cue, and whether explicit KV-related signals provide complementary information near the regime boundary.

\subsection{Claims and Evaluation Logic}

Our experimental logic follows three questions.
\begin{enumerate}
\item \textbf{Regime characterization}: when does realistic dynamic offload show positive value, and when does it fail?
\item \textbf{Predictive signals}: is the transition more consistent with context alone, attention share, or KV-related pressure signals?
\item \textbf{Decision implication}: do static attention rules and generic realistic heuristics over- or under-predict early-positive cases, and can KV-aware features refine the boundary?
\end{enumerate}

These questions motivate the main figures and tables in the paper. \Cref{fig:context-transition} addresses the first question, \Cref{fig:kv-pressure} and \Cref{fig:decision-signal-study} address the second, and \Cref{fig:policy-comparison} provides a focused case-study view for the third.
